\chuong{GIỚI THIỆU TỔNG QUAN}

\subsection{Bối cảnh thực tế}

Dữ liệu bán hàng được ghi nhận qua từng đơn hàng và dòng sản phẩm, còn người quản lý cần nhìn thấy doanh thu thay đổi ra sao theo thời gian, sản phẩm và khu vực. Để trả lời những câu hỏi đó, dữ liệu giao dịch phải được tổng hợp và đặt trong một cấu trúc có thể tra cứu theo nhiều góc nhìn.  Chaudhuri và Dayal~\citeyearpar{chaudhuri1997} cho thấy kho dữ liệu phục vụ phân tích cần dữ liệu lịch sử, được tích hợp và tổ chức phù hợp với các truy vấn tổng hợp.

Adventure Works Cycles là doanh nghiệp giả định trong một tình huống hướng dẫn của Microsoft. Trong tình huống này, dữ liệu giao dịch và thông tin ngoài giao dịch như hạn mức doanh số từ bảng tính đã được tập hợp vào kho dữ liệu quan hệ AdventureWorksDW2019. Microsoft mô tả các hạn chế ở tầng khai thác: báo cáo tĩnh khó cho phép đi sâu vào dữ liệu, người dùng mất thời gian tạo truy vấn do cấu trúc phức tạp và cách tính được viết riêng cho từng báo cáo có thể làm số liệu tổng hợp khác nhau. Tài liệu đánh giá kho dữ liệu này được thiết kế tốt và không có vấn đề lớn ở khâu nạp dữ liệu. Microsoft lựa chọn bổ sung Analysis Services để cải thiện truy vấn và tái sử dụng cách tính~\citep{microsoft-scenario}.

Tình huống hướng dẫn bắt đầu từ một kho dữ liệu đã có và tập trung khai thác dữ liệu bằng Analysis Services. Đồ án này sử dụng Adventure Works làm bối cảnh để thực hành thiết kế kho dữ liệu từ các bảng giao dịch của bộ dữ liệu. Công việc bao gồm xác định mức chi tiết của bảng sự kiện, thiết kế các chiều phân tích và thống nhất quy tắc tính chỉ tiêu, nhằm giải quyết hai hạn chế mà tình huống mẫu nêu ra là khó truy vấn theo nhiều góc nhìn và số liệu khác nhau giữa các báo cáo, trước khi kiểm tra kết quả trên báo cáo.

Microsoft Fabric được chọn làm nền tảng triển khai vì tài liệu tổng quan của Microsoft mô tả khả năng nạp, chuyển đổi, lưu trữ, phân tích và báo cáo trong cùng môi trường~\citep{microsoft-fabric}. Trong đó, mục Warehouse cho phép tổ chức dữ liệu theo cấu trúc phân tích; Microsoft~\citeyearpar{microsoft-warehouse} nêu mô hình hình sao và hình tuyết là các trường hợp phù hợp với Fabric Data Warehouse. Trên nền tảng này, đồ án xây dựng kho dữ liệu, quy trình cập nhật, báo cáo Power BI, Data Agent và các mô hình học máy hỗ trợ dự báo lượng bán và hành vi khách hàng.

\subsection{Lý do chọn đề tài}

Bài toán bán hàng của Adventure Works Cycles được chọn vì dữ liệu có thể liên kết giao dịch với sản phẩm, khách hàng, thời gian và khu vực. Các mối liên kết này cho phép thiết kế bảng sự kiện, bảng chiều và kiểm tra những chỉ tiêu cụ thể như doanh thu, số lượng bán và số đơn hàng. Nhờ vậy, số liệu trên báo cáo có thể được đối chiếu ngược với dữ liệu nguồn.

Để các chỉ tiêu được tính nhất quán, dữ liệu cần được chuẩn hóa trong một quy trình có thể chạy lại. Fabric cho phép triển khai quy trình đó cùng với kho dữ liệu và lớp báo cáo trong một môi trường. Dữ liệu bán hàng đã xử lý trong kho đồng thời là đầu vào cho các bài toán dự báo lượng bán, khả năng mua lại và nhóm sản phẩm khách hàng có thể mua tiếp.

Trong phạm vi đồ án môn học, mỗi quyết định thiết kế có thể được kiểm tra bằng kết quả cụ thể: dữ liệu sau xử lý có khớp nguồn hay không, các chỉ tiêu có cho cùng kết quả khi đối chiếu, và quy trình cập nhật có chạy ổn định. Những kết quả này cũng cho thấy các khâu cần điều chỉnh.

\subsection{Mục tiêu dự án}

\subsubsection{Mục tiêu tổng quan}

Mục tiêu tổng quan của dự án là xây dựng giải pháp kho dữ liệu trên nền tảng Microsoft Fabric phục vụ việc tích hợp, quản lý và khai thác dữ liệu bán hàng của Adventure Works Cycles. Giải pháp hướng đến việc tổ chức dữ liệu thành một hệ thống tập trung, nhất quán và thuận tiện cho phân tích đa chiều, từ đó hỗ trợ theo dõi tình hình kinh doanh, đánh giá hiệu quả bán hàng và nhận diện các xu hướng từ dữ liệu lịch sử.

Trên cơ sở kho dữ liệu được xây dựng, dự án kết hợp ứng dụng học máy, các quy trình xử lý dữ liệu theo từng chặng và trực quan hóa kết quả phân tích nhằm hỗ trợ ra quyết định trong hoạt động bán hàng. Đồng thời, thông qua quá trình triển khai, thử nghiệm và đánh giá, dự án làm rõ khả năng ứng dụng nền tảng Microsoft Fabric trong việc xây dựng một giải pháp dữ liệu từ tích hợp dữ liệu nguồn đến khai thác thông tin phục vụ nhu cầu kinh doanh.

\subsubsection{Mục tiêu cụ thể}

Để đạt được mục tiêu tổng quan, dự án hướng đến các mục tiêu cụ thể sau:

Thứ nhất, khảo sát, mô tả và phân tích dữ liệu bán hàng của Adventure Works Cycles nhằm xác định cấu trúc dữ liệu, mối quan hệ giữa các bảng và các thông tin cần thiết cho bài toán phân tích. Đồng thời, đánh giá chất lượng dữ liệu, nhận diện các vấn đề về dữ liệu thiếu, trùng lặp hoặc không nhất quán để làm cơ sở xây dựng quy trình xử lý phù hợp.

Thứ hai, xây dựng quy trình tích hợp và xử lý dữ liệu trên nền tảng Microsoft Fabric, bao gồm thu thập, nạp, làm sạch, chuẩn hóa và chuyển đổi dữ liệu. Quy trình này hướng đến việc bảo đảm tính đầy đủ, nhất quán và toàn vẹn của dữ liệu trước khi đưa vào kho dữ liệu phục vụ phân tích và khai thác.

Thứ ba, thiết kế và triển khai kho dữ liệu theo mô hình đa chiều, tổ chức các bảng sự kiện và bảng chiều phù hợp với đặc điểm dữ liệu bán hàng. Kho dữ liệu cần đáp ứng nhu cầu phân tích các chỉ tiêu như doanh thu, số lượng sản phẩm bán ra và số lượng đơn hàng theo những chiều phân tích chủ yếu, bao gồm thời gian, sản phẩm, khách hàng và khu vực bán hàng.

Thứ tư, xây dựng và đánh giá ba mô hình học máy phục vụ hoạt động bán hàng trên cơ sở dữ liệu đã được xử lý và tổ chức trong kho dữ liệu. Ba bài toán lần lượt dự báo lượng bán của sản phẩm trong 7 ngày, khả năng khách hàng mua lại trong 30 ngày và nhóm sản phẩm khách hàng có thể mua tiếp trong 60 ngày. Kết quả cung cấp thêm thông tin tham khảo cho việc chuẩn bị hàng hóa và chăm sóc khách hàng.

Thứ năm, xây dựng và kiểm tra các quy trình xử lý dữ liệu cho từng chặng, gồm nạp dữ liệu vào Bronze, chuyển đổi sang Silver, công bố dữ liệu Gold và chuẩn bị dữ liệu học máy. Kết quả triển khai là cơ sở đánh giá khả năng hợp nhất các quy trình này thành một luồng điều phối đầu cuối trong giai đoạn tiếp theo.

Thứ sáu, xây dựng các báo cáo và dashboard trực quan phục vụ theo dõi, phân tích và đánh giá tình hình bán hàng. Các báo cáo hướng đến việc thể hiện những chỉ tiêu kinh doanh chủ yếu, hỗ trợ so sánh kết quả giữa các thời kỳ, sản phẩm và khu vực, đồng thời trình bày kết quả dự báo để người sử dụng thuận tiện khai thác thông tin.

Thứ bảy, thử nghiệm và đánh giá giải pháp dựa trên các tiêu chí về chất lượng dữ liệu, tính chính xác của kết quả tổng hợp, hiệu quả truy vấn, khả năng thực thi của từng quy trình dữ liệu và mức độ đáp ứng của mô hình học máy. Trên cơ sở kết quả thực nghiệm, xác định những hạn chế của giải pháp và đề xuất hướng hoàn thiện, mở rộng trong tương lai.

\subsection{Đối tượng và phạm vi dự án}

\subsubsection{Đối tượng nghiên cứu}

Đối tượng nghiên cứu của dự án là dữ liệu bán hàng của Adventure Works Cycles và quy trình xây dựng giải pháp kho dữ liệu trên nền tảng Microsoft Fabric phục vụ khai thác dữ liệu này. Dự án tập trung vào các thông tin liên quan đến đơn hàng, chi tiết bán hàng, sản phẩm, khách hàng, thời gian và khu vực bán hàng, cùng với mối quan hệ giữa các đối tượng dữ liệu trong quá trình phân tích.

Dự án nghiên cứu mô hình tổ chức dữ liệu trong kho dữ liệu, các quy trình tích hợp và chuyển đổi theo từng chặng, cùng phương pháp ứng dụng học máy cho bài toán bán hàng. Các báo cáo và dashboard được xây dựng hướng đến nhu cầu sử dụng của bộ phận quản lý và phân tích bán hàng, phục vụ theo dõi kết quả kinh doanh và tham khảo thông tin dự báo.

\subsubsection{Phạm vi nghiệp vụ bán hàng}
Dự án tập trung vào đơn hàng, chi tiết đơn hàng, sản phẩm, khách hàng và khu vực bán hàng. Các chỉ tiêu được xem xét gồm doanh thu, số lượng bán và số đơn hàng theo thời gian, sản phẩm, khách hàng và khu vực.

\subsubsection{Phạm vi dữ liệu}

Về phạm vi dữ liệu, dự án sử dụng dữ liệu Adventure Works được nhóm lựa chọn để triển khai, tập trung vào các bảng liên quan đến hoạt động bán hàng và những bảng nghiệp vụ cần thiết để bổ sung thông tin phân tích. Phạm vi thời gian được giới hạn trong khoảng dữ liệu lịch sử có sẵn của bộ dữ liệu sử dụng. Việc phân tích và đánh giá được thực hiện trên bộ dữ liệu này trong bối cảnh doanh nghiệp giả định.

\subsubsection{Phạm vi triển khai}

Về phạm vi xây dựng hệ thống, dự án triển khai các quy trình thu thập, nạp, làm sạch, chuẩn hóa và chuyển đổi dữ liệu trên nền tảng Microsoft Fabric. Dữ liệu sau xử lý được tổ chức thành kho dữ liệu theo mô hình đa chiều, gồm các bảng sự kiện và bảng chiều phù hợp với yêu cầu phân tích bán hàng. Mỗi chặng được triển khai và kiểm tra riêng; một pipeline điều phối thống nhất cho toàn bộ luồng chưa thuộc kết quả hoàn thành của phiên bản hiện tại.

Về phạm vi phân tích, dự án tập trung vào các chỉ tiêu chủ yếu như doanh thu, số lượng sản phẩm bán ra và số lượng đơn hàng. Các chỉ tiêu được khai thác theo thời gian, sản phẩm, khách hàng và khu vực bán hàng nhằm theo dõi xu hướng, so sánh kết quả và nhận diện những đặc điểm đáng chú ý trong dữ liệu. Kết quả phân tích được thể hiện thông qua các báo cáo và dashboard trực quan.

Về phạm vi học máy, dự án triển khai ba bài toán ở hai mức hạt. Bài toán thứ nhất dự báo lượng bán 7 ngày tới cho từng sản phẩm tại từng ngày chốt. Hai bài toán còn lại được xây dựng theo khách hàng và ngày chốt, nhằm dự đoán khả năng mua lại trong 30 ngày và nhóm sản phẩm của lần mua tiếp theo trong 60 ngày. Các tập huấn luyện và kiểm tra được tách theo thời gian để phản ánh tình huống dự báo trên dữ liệu tương lai.

Về phạm vi thử nghiệm và đánh giá, dự án kiểm tra chất lượng dữ liệu sau xử lý, tính chính xác của các chỉ tiêu tổng hợp, khả năng truy vấn và khả năng thực thi của từng quy trình dữ liệu đơn lẻ. Đối với mô hình học máy, việc đánh giá được thực hiện trên dữ liệu kiểm tra theo trình tự thời gian để xem xét mức độ phù hợp của kết quả dự báo. Kết quả thử nghiệm là cơ sở xác định những hạn chế và đề xuất hướng phát triển của giải pháp.

Dự án được triển khai ở quy mô đồ án môn học, tập trung mô phỏng và đánh giá giải pháp trong môi trường thử nghiệm. Các nội dung như vận hành hệ thống ở quy mô doanh nghiệp, tích hợp dữ liệu giao dịch thời gian thực, tối ưu toàn bộ hoạt động sản xuất và quản lý chuỗi cung ứng nằm ngoài phạm vi triển khai của đồ án.

\subsection{Công nghệ và nền tảng sử dụng}

Kiến trúc của dự án được thiết kế với nguồn dữ liệu Azure SQL Database, các tầng xử lý trong Microsoft Fabric và lớp khai thác dữ liệu. Mỗi công nghệ được chọn cho một phần của luồng dữ liệu, từ tiếp nhận dữ liệu giao dịch đến lưu trữ, mô hình hóa, học máy và báo cáo.

\subsubsection{Azure SQL Database}

Azure SQL Database được dùng để mô phỏng hệ thống giao dịch của Adventure Works Cycles. Dữ liệu CSV của bộ AdventureWorks được nạp vào cơ sở dữ liệu và giữ cấu trúc theo năm lược đồ nghiệp vụ: HumanResources, Person, Production, Purchasing và Sales. Trong kiến trúc, Azure SQL Database là nguồn OLTP; dữ liệu từ đây được sao chép vào Fabric để xử lý mà không dùng trực tiếp cho các truy vấn phân tích.

\subsubsection{Microsoft Fabric và OneLake}

Microsoft Fabric là môi trường được chọn để triển khai các thành phần tích hợp, lưu trữ, xử lý và khai thác dữ liệu của dự án~\citep{microsoft-fabric}. Theo thiết kế, các Lakehouse, Warehouse, pipeline và công cụ học máy được tổ chức trong một workspace. Semantic model nối dữ liệu tầng Gold với lớp khai thác. Bên dưới các đối tượng Fabric, OneLake cung cấp lớp lưu trữ dùng chung, cho phép các thành phần truy cập dữ liệu trong một không gian lưu trữ thống nhất~\citep{microsoft-onelake}.

\subsubsection{Lakehouse và Medallion Architecture}

Lakehouse được bố trí tại các tầng Bronze, Silver và nhánh học máy của tầng Gold. Tầng Bronze sẽ lưu các bảng nguồn ở dạng thô để phục vụ đối chiếu và chạy lại quy trình. Tầng Silver gồm vùng STG để chuẩn hóa dữ liệu và vùng INT để tích hợp các bảng, áp dụng quy tắc nghiệp vụ. Tại tầng Gold, một Lakehouse riêng được dành cho tập đặc trưng, biến mục tiêu và kết quả dự báo. Cách phân lớp này bám theo Medallion Architecture, trong đó dữ liệu được nâng dần về chất lượng và mức độ sẵn sàng qua từng tầng~\citep{microsoft-medallion}.

\subsubsection{Fabric Data Factory}

Fabric Data Factory được chọn để di chuyển và điều phối dữ liệu giữa các thành phần của kiến trúc. Theo luồng thiết kế, Data Pipeline sẽ sao chép dữ liệu từ Azure SQL Database vào Lakehouse tầng Bronze, gọi Dataflow Gen2 để xử lý dữ liệu từ Bronze sang Silver và thực thi các script nạp dữ liệu từ Silver vào Warehouse tầng Gold. Việc gom các bước này vào pipeline xác định thứ tự phụ thuộc và tạo cơ sở cho việc lập lịch, theo dõi trạng thái chạy và xử lý lỗi.

\subsubsection{Fabric Data Warehouse và Semantic Model}

Fabric Data Warehouse được bố trí tại tầng Gold để lưu dữ liệu đã tổ chức theo mô hình chiều. Warehouse chứa các bảng chiều, bảng sự kiện và data mart phục vụ phân tích bán hàng; đây cũng là nguồn của Semantic Model~\citep{microsoft-warehouse}. Semantic Model định nghĩa các mối quan hệ và chỉ tiêu dùng chung trước khi dữ liệu được Power BI hoặc Data Agent truy vấn. Cách tổ chức này tách phần xử lý, lưu trữ dữ liệu khỏi phần trình bày báo cáo.

\subsubsection{Fabric Data Science và Machine Learning}

Nhánh học máy nhận các bảng đặc trưng từ Lakehouse tầng Gold. Ba notebook huấn luyện mô hình cố định, theo dõi thí nghiệm bằng MLflow và đăng ký ba ML Model tương ứng trong workspace. Kết quả chấm điểm được ghi trở lại Lakehouse tầng Gold để kết hợp với dữ liệu lịch sử khi phân tích. Biến đầu vào, biến mục tiêu, phương pháp tách tập dữ liệu và chỉ số đánh giá sẽ được trình bày tại chương xây dựng mô hình học máy.

\subsubsection{Power BI và Data Agent}

Power BI và Data Agent được xác định là hai đối tượng khai thác dữ liệu ở cuối kiến trúc. Power BI sẽ sử dụng Semantic model để xây dựng dashboard theo dõi doanh thu, số lượng bán, số đơn hàng và các chiều phân tích. Data Agent sẽ khai thác cùng mô hình để hỗ trợ truy vấn dữ liệu theo ngữ cảnh nghiệp vụ. Việc dùng chung lớp ngữ nghĩa giúp các quan hệ và cách tính chỉ tiêu không phải định nghĩa lại cho từng công cụ~\citep{microsoft-fabric}.
